\documentclass[10pt,a4paper]{article}
\usepackage[T1]{fontenc}
\usepackage[margin=22mm]{geometry}
\usepackage{lmodern,amsmath,amssymb,amsthm,mathtools,booktabs,microtype}
\usepackage[hidelinks]{hyperref}
\newtheorem{theorem}{Theorem}
\newtheorem{lemma}[theorem]{Lemma}
\title{A paired-divisor lower bound for a dyadic Erd\H{o}s--Straus selector}
\author{The Clankers}
\date{15 September 2026 --- Research preprint}
\begin{document}\maketitle\vspace{-3mm}
\begin{abstract}
An actual bounded block of powers of three gives a pointwise lower bound for
middle-channel states with a fixed divisor seed. We prove its exact count,
constructive return, and sharp multiplicity threshold. Complementary-factor
reflection retains an explicit availability defect. No universal
Erd\H{o}s--Straus or historical-priority claim is made.
\end{abstract}

\paragraph{Original arithmetic and return.}
Let $p\equiv1\pmod4$ be an integer, $u\ge1$, $p>2u$ and $\gcd(p,u)=1$.
Put $K=\prod_{\ell^f\parallel u}\ell^{\lceil f/2\rceil}$, $m=4K$, $N=p+4u$.
An original middle state with fixed $u$ and $p/4<a<p/2$ is exactly an ordered
factor pair
\[
 \boxed{RQ=N,\qquad Q\equiv-1\pmod m.}\tag{1}
\]
Indeed $u\mid a^2$ iff $K\mid a$; the middle gate is $R\mid p+4u$ and
$R=4a-p$. These imply (1), since $N\equiv p\pmod m$.
Conversely (1) implies $R\equiv-p\pmod m$, $R,Q\equiv3\pmod4$,
$3\le R,Q<p$, and $K\mid a=(p+R)/4$. Set
\[
 d=\gcd(a,u),\quad(h,r,s)=(d^2/u,u/d,a/d),\quad\lambda=(r+s)/R.
\]
Prime valuations give $a=hrs$, $u=hr^2$, $\gcd(r,s)=1$. Since $hr$ is a unit
modulo $R$, the middle gate makes $\lambda$ a positive integer and
\[
 \boxed{\frac4p=\frac1{hrs}+\frac1{phs\lambda}+\frac1{phr\lambda}.}\tag{2}
\]
Clearing reciprocals proves (2). For the ordered triple $(a,y,z)$ its divisor
inverse is $u=pa^2/(Ry-pa)$. Thus all original exponents and the order are
recovered, not merely a rational point. This is the inherited Type II
parametrization \cite[\S2, (2.13)--(2.22), Propositions 2.6--2.7]{ET}.

\paragraph{An exact same-parameter reflection obstruction.}
The swap $(R,Q)\mapsto(Q,R)$ is an involution of the ordered factor pairs.
Its proposed new shell $a_Q=(p+Q)/4$ still has a zero middle gate, but
\[
 \boxed{
 \frac K{\gcd(K,a_Q)}=\frac K{\gcd(K,(p-1)/4)},\qquad
 \frac u{\gcd(u,a_Q^2)}=\frac u{\gcd(u,((p-1)/4)^2)}.}\tag{3}
\]
For the first equality use $a_Q\equiv(p-1)/4\pmod K$.
For $\ell^f\parallel u$, a valuation below $\lceil f/2\rceil$ agrees on both
sides; otherwise both squared valuations reach $f$. This proves the second.
An existing state therefore reflects to another original state iff
$p\equiv1\pmod{4K}$. On that locus the smaller of $R,Q$ gives a residual
at most $\sqrt N$. Outside it a zero gate does not supply a missing divisor.

\begin{theorem}[Dyadic mass and occupancy]\label{thm:main}
Let $k\ge4$ and
\[
 m=2^k,\quad o=2^{k-2},\quad L=o/2,\quad u=2^{2k-5},
 \qquad p>2u,\quad p\equiv1+2^{k-1}\pmod m.
\]
Write $N=p+4u=3^eM$ with $3\nmid M$, and define
\[
 A=\#\{t\mid M:t\equiv5,7\pmod8\}
 =\frac12\left(\prod_{q^f\parallel M}(f+1)
 -\prod_{q^f\parallel M}\sum_{j=0}^f\chi(q)^j\right),\tag{4}
\]
where $\chi=1$ at odd residues $1,3\pmod8$ and $-1$ at $5,7\pmod8$.
If $C$ counts the original middle states (1), then
\[
 \boxed{
 \frac12\left\lfloor\frac{2(e+1)}o\right\rfloor A
 \le C\le
 \frac12\left\lceil\frac{2(e+1)}o\right\rceil A.}\tag{5}
\]
When $L\mid e+1$, $C=(e+1)A/o$. In particular, for $e\ge L-1$,
\[
 \boxed{C>0\iff\text{some prime }q\mid N\text{ is }5\text{ or }7\pmod8.}\tag{6}
\]
\end{theorem}

\newpage
\begin{proof}
Induction gives $v_2(3^{2^t}-1)=t+2$ for $t\ge1$: factor the next difference
as $(3^{2^t}-1)(3^{2^t}+1)$, whose second factor has valuation one.
Thus $3$ has order $o$ modulo $m$ and $3^L=1+2^{k-1}\pmod m$.
Its subgroup $H$ consists exactly of units $1,3\pmod8$: it is contained there
and the sizes agree. All other units lie in $-H$, and $-1\notin H$.
We have $N\equiv p\equiv3^L\pmod m$.

Let $a_j=\#\{t\mid M:t\equiv-3^j\pmod m\}$, $0\le j<o$.
The original residual target is $-p=-3^L$. Unique factorization gives the
bijection
\[
 (t,i)\mapsto R=3^it,\qquad t\mid M,\quad0\le i\le e,\quad i+j=L\pmod o,
\]
with inverse $i=v_3(R)$ and $t=R/3^i$. Therefore
\[
 C=\sum_j a_jn_j,\qquad n_j=\#\{0\le i\le e:i+j=L\pmod o\}.
\]
The actual divisor involution $t\mapsto M/t$ has index
$j'=L-e-j\pmod o$, so $a_j=a_{j'}$.
Write $e+1=fo+s$, $0\le s<o$. Then $n_j=f+\mathbf1_T(j)$ for the cyclic
interval $T=[L-s+1,L]$. The involution carries $T$ to $T-L=[1-s,0]$.
These intervals are disjoint for $s<L$, partition the cycle for $s=L$, and
cover it for $s>L$. Hence
\[
 \left\lfloor2(e+1)/o\right\rfloor\le n_j+n_{j'}
 \le\left\lceil2(e+1)/o\right\rceil.
\]
Averaging the unchanged arithmetic mass,
$C=\frac12\sum_j a_j(n_j+n_{j'})$, proves (5), including fixed points of the
involution. The bounds agree exactly at $s=0,L$, proving the stated equality.
Formula (4) is the expansion of the multiplicative character on the actual
divisor set. Its mass is positive iff an outside prime occurs. Finally
$e\ge L-1$ makes the lower bound at least $A/2$, proving (6).
\end{proof}

\paragraph{Constructive return from one actual prime.}
For an outside prime $q$, let $q\equiv-3^j\pmod m$ with $0\le j<o$.
Of $i_Q=-j\pmod o$ and $i_R=L-j\pmod o$, exactly one has representative
in $[0,L-1]$. The threshold makes that exponent available among the actual
factors $3^e$. If it is $i_Q$, choose $Q=3^{i_Q}q$, $R=N/Q$; otherwise choose
$R=3^{i_R}q$, $Q=N/R$. This satisfies (1) and returns (2).
The prime, exponent and residual/cofactor role remain in the certificate.
The logarithm can be found by an exact lift from modulo eight: at level $r$
keep $j$ or add $2^{r-3}$, as the two powers differ by $1+2^{r-1}$ modulo
$2^r$. Both the algorithm and its inverse check are finite.

All such reflected states have the two defects in (3) equal to two, because
$v_2((p-1)/4)=k-3$ and $v_2(u)=2k-5$. This does not obstruct the direct
construction just proved. If the stronger actual capacity $e\ge2o-2$ holds,
choose the least outside prime $q$. Outside prime occurrences have even total
multiplicity since $\chi(M)=1$, so $q^2\le M$. The residual choice with
$i_R\le o-1\le e/2$ gives $R\le\sqrt N$ and $\Omega(R)\le o$.

\paragraph{Sharpness of the exponent threshold.}
For $k\ge5$ choose a sufficiently large prime $q\equiv-3\pmod{2^k}$ and set
\[
 e=L-2,\quad N=3^{L-2}q^2,\quad p=N-4u.
\]
Then all parameter conditions hold and $A=1$. Every outside divisor is
$3^iq$, $0\le i\le L-2$, whereas the target needs $i+1=L\pmod o$.
Thus $C=0$. Such $q$ exist by Dirichlet's theorem \cite[Theorem 18.1]{Dir};
this does not claim primality of every derived $p$.
Actual hard-prime instances are
\[
\begin{array}{c|r|r|r|r}
k&u&e&q&p\\\hline
5&32&2&829&6185041\\
6&128&6&3517&9017211169.
\end{array}
\]
The complete fixed-seed sets are empty in both cases despite an outside prime.
Their primality is deterministically checked, and other ES witnesses are
supplied. No counterexample to ES is asserted.

\newpage
\paragraph{An explicit same-prime branch with an unavoidable long residual.}
At $k=4$ the seed is eight. For $p\equiv25\pmod{48}$ its required capacity
$e\ge1$ is automatic, so (6) is an exact criterion for this entire progression.
The return for $q\bmod16=7,13,15,5$ is respectively
\[
 R=q,\quad3q,\quad N/q,\quad N/(3q).
\]
At the certified hard prime $p=2542201$,
\[
 N=p+32=3\cdot37^2\cdot619,\qquad A=2,\quad C=A/2=1.
\]
Its unique seed-eight residual is $R=37\cdot619=22903$ and $Q=111$.
Since $R^2>N$, this branch cannot supply a square-root residual merely by
swapping factors. The original data are
\[
 (a,h,r,s,\lambda)=(641276,2,2,160319,7),
\]
\[
 \boxed{\frac4{2542201}
 =\frac1{641276}+\frac1{5705883709666}+\frac1{71181628}.}
\]
Its reflected shell $a_Q=635578$ has $8\nmid a_Q^2$; exactly one factor two
is missing. The reflected gate remains zero, demonstrating why availability
must accompany the cofactor map.

\paragraph{Large multiplicities without expanding every word.}
Formula (4) uses all prime multiplicities but only a product of local integer
sums. Formula (5) therefore controls an unbounded exponent array by the same
explicit constants. The full count outside the equality cases remains the
finite cyclic convolution $\sum a_jn_j$, with every coefficient's original
divisor fibre retained. No event is assumed before obtaining the lower bound.
Factoring $N$ is still required for this implementation; no polynomial-time
factorization or universal prime coverage is claimed.

\paragraph{Relation to the SplitZero source.}
The original carrier distinguishes absence $\tau$ from supported zero
$e=0^\bullet$ \cite{Previous}. In (1)--(3), a reflected gate may evaluate to
$e$ while its divisor is not available in the original exponent box. The
proposal, its zero gate, and the missing prime power are separate retained
records. In the count (4), coefficients instead count actual divisor words;
there is no signed cancellation. The factor one-half in the proof averages a
bijection of the original divisor set, not independent local candidates.
No additional cohomological identification is needed to conclude (6).

\paragraph{Reproducibility and scope.}
The accompanying standard-library verifier and an independent divisor-enumeration
checker test the fixed-seed correspondence, original ordered inverses, reflection
valuations, exact cyclic counts, and the sharpness examples. Over the 2245 primes
$p\le2000000$ in $25\pmod{48}$ and the six hard classes modulo 840, there are
839 occupied seed-eight branches and 1303 original middle states. In 240 occupied
branches the minimum residual exceeds $\sqrt{p+32}$. These are fixed-seed
certificates, not a new ES verification range. The general proofs do not depend
on that census.

A universal proof would still have to force one usable seed branch for each
prescribed prime. The exact relations between the inputs,
$(p+4v)-(p+4u)=4(v-u)$, do not by themselves exclude simultaneous failure.
This note proves a within-branch arithmetic lower bound and a sharp capacity
criterion, not that additional cross-branch implication.

\begingroup\small
\begin{thebibliography}{3}
\bibitem{ET} Elsholtz, C., \& Tao, T. (2013). Counting the number of solutions
to the Erd\H{o}s--Straus equation on unit fractions. \emph{Journal of the Australian
Mathematical Society, 94}(1), 50--105. arXiv:1107.1010v6, \S2.
\bibitem{Dir} Sutherland, A. V. (2017). \emph{Dirichlet L-functions, primes in
arithmetic progressions}. MIT 18.785 Lecture 18, Theorem 18.1.
\bibitem{Previous} The Clankers (2026). \emph{Split scalar arithmetic and the
residual-27 catalogue}. Attached predecessor, September 15. The source-reading
ledger in the full package records its hashes, the earlier fixed-seed CRT
locators, the content search, and the inherited/new statement boundary.
\end{thebibliography}
\endgroup
\end{document}
